
\begin{thm}[{Edge universality on generalized Wigner matrix: \cite[Theorem 2.7 and Proposition 6.4, (6.15) at $f=0$]{EKYY2}}] \label{thm: old edge un}
Let $H^{\w}$ be a GOE and $H^{\V}$ a generalized symmetric Wigner matrix with
$$
\E^{\V} |h_{ij}|^2 = N^{-1} + \delta_{ij} N^{-1}.
$$
Assume that $H^{\V}$ satisfies the bounded support condition with $q=N^{\phi}$, for some constant $1/3 < \phi \leq 1/2$. Then, there exist constants $\delta > 0$ and $N_0$ such that, for any $N \geq N_0$ and any $s \in \mathbb{R}$, we have
\begin{equation} \label{qset}
\p^{\w} \Big( N^{2/3} (\lambda_N -2) \leq s - N^{-\delta} \Big) - N^{-\delta} \leq \p^{\V} \big( N^{2/3} (\lambda_N -2) \leq s \big) \leq \p^{\w} \Big( N^{2/3} (\lambda_N -2) \leq s + N^{-\delta} \Big) + N^{-\delta}.
\end{equation}
Here, $\p^{\V}$ and $\p^{\w}$ denote the laws of the ensembles $H^{\V}$ and $H^{\w}$, respectively, and $\delta$ and $N_0$ do not depend on $s$.
\end{thm}

The diagonal variance here and for the GOE is $2/N$, while it is $1/N$ in \cite[Definition 2.1]{EKYY2}. For entries bounded by $N^{-\phi}$ almost surely, Theorem \ref{thm: old edge un} is the case $C_0=1$ of Lemma \ref{lem: ex edge un} below. In the form stated above, with the bounded support condition of Definition \ref{cond:bdd}, it is cited from \cite{EKYY2} and is not used separately in this paper.

\begin{rem} \label{genedge} As in \cite{EYY} and \cite{EKYY2}, Theorem \ref{thm: old edge un}, as well as Lemma \ref{lem: ex edge un} and Theorem \ref{thm: new edge un} below, can be extended to finite correlation functions of extreme eigenvalues. For example, we have the following extension:
 \begin{align}\label{twa}
&  \p^{\bf w} \Big ( N^{2/3}  ( \lambda_N -2) \le s_1- N^{-\e}, \ldots, N^{2/3} ( \lambda_{N-k} -2)
\le s_{k+1}- N^{-\e} \Big )- N^{-\delta}   \nonumber \\
&
 \le   \p^{\bf v} \Big ( N^{2/3} (  \lambda_N -2) \le s_1,  \ldots, N^{2/3} ( \lambda_{N-k} -2)
\le s_{k+1}  \Big )  \\
 &  \le
 \p^{\bf w} \Big ( N^{2/3} ( \lambda_N -2) \le s_1+ N^{-\e}, \ldots, N^{2/3} ( \lambda_{N-k} -2)
 \le s_{k+1}+ N^{-\e}   \Big )+ N^{-\delta}  \nonumber
\end{align}
for all $k$ fixed and $N$ sufficiently large. The proof of \eqref{twa} is similar to that of \eqref{qset} except that it uses the general form of the Green function comparison theorem. This remark is not used in this paper (see \eqref{segg}), and it is not formalized.
\end{rem}

We slightly extend this result as follows:

\begin{lem}\label{lem: ex edge un}
Let $H^{\w}$ be a GOE and $H^{\V}$ a generalized symmetric Wigner matrix satisfying \eqref{eq:diag} with a constant $C_0\ge0$, whose entries satisfy $|h_{ij}|\le N^{-\phi}$ almost surely for some $1/3<\phi\le1/2$. Then there are $\delta>0$ and $N_0$, depending only on $(C_0,\phi)$, such that \eqref{qset} holds for all $N\ge N_0$ and all $s\in\mathbb R$.
\end{lem}

This is the almost sure form of the bounded support condition, the only case used here (Theorem \ref{thm: new edge un} applies it to $\wt H$ of Lemma \ref{lem: match}); the probabilistic form is not proved here. We postpone the proof of this lemma to the end of this section.

\bigskip

To prove our main result, we claim the following three lemmas, which extend the previous results to the Wigner matrix with bounded support condition of small $q$. First, we improve the previous result on rigidity. We define the {\it normalized empirical counting function} by
$$
{\mathfrak n}(E):= \frac{1}{N}\# \{ \lambda_j\le E\}.
$$
Let
$$
n_{sc}(E) :  = \int_{-\infty}^E \varrho_{sc}(x) {\rm d} x
$$
be the distribution function of the semicircle law.

\begin{thm}[Rigidity of eigenvalues: small $q$ case] \label{thm: new rigidity}
Let $H$ be a generalized symmetric Wigner matrix with some constant $C_1$ such that for any $1 \le i< j\leq N$
$$
\E |h_{ii}|^2 \leq C_1/N, \quad \E  |h_{ij}|^2=1/N, \quad \E |h_{ij}|^3 \leq C_1 N^{-3/2}, \quad \E  |h_{ij}|^4 \leq C_1(\log N) N^{-2}
$$
and $|h_{ij}|\le N^{-\phi}$ almost surely for all $1\le i,j\le N$, for some constant $\phi>0$. Then, there exist constants $C>0$ and $N_0$, depending only on $C_1$ and $\phi$, such that with high probability we have
\be \label{rigiditynew}
\bigcap_j \bigg\{ |\lambda_j-\gamma_j| {\,\leq\,} \varphi^C \Big [ \min \big ( \, j ,  N-j+1 \, \big) \Big]^{-1/3} N^{-2/3} \bigg\}
\ee
and
\be\label{rigiditynew2}
\sup_{E\in\mathbb R} \big| {\mathfrak n} (E)-n_{sc}(E)\big|  {\,\leq\,}  \frac{\varphi^C}{N}
\ee
for any $N > N_0$.
\end{thm}

Theorem \ref{thm: new rigidity} is the special case, with entries bounded almost surely, of the corresponding statement under the bounded support condition of Definition \ref{cond:bdd} with $q=N^{\phi}$. That statement is not proved here: under Definition \ref{cond:bdd} alone the moment bound (6.78) can fail. Section \ref{sec4} uses only the almost sure case, since $\wt H^S$ is bounded almost surely.

Next theorem shows that the edge universality holds under the assumptions in Theorem \ref{thm: new rigidity}.

\begin{thm}[Edge universality: small $q$ case] \label{thm: new edge un}
Let $H^{\w}$ be a GOE and $H^\V$ be a generalized symmetric Wigner matrix satisfying the conditions for $H$ in Theorem \ref{thm: new rigidity} (in particular $|h_{ij}|\le N^{-\phi}$ almost surely). Then, there exist constants $\delta>0$ and $N_0$, depending only on $C_1$ and $\phi$, such that for all $N\ge N_0$, all such $H^{\V}$ and all $s\in \mathbb R$, we have
\begin{equation} \label{qset1}
\p^{\w} \Big( N^{2/3} (\lambda_N -2) \leq s - N^{-\delta} \Big) - N^{-\delta} \leq \p^{\V} \big( N^{2/3} (\lambda_N -2) \leq s \big) \leq \p^{\w} \Big( N^{2/3} (\lambda_N -2) \leq s + N^{-\delta} \Big) + N^{-\delta}.
\end{equation}
Here, $\p^{\V}$ and $\p^{\w}$ denote the laws of the ensembles $H^{\V}$ and $H^{\w}$, respectively.
\end{thm}

Finally, we show a weak bound on $G_{ij}$ ($i\neq j$) of $H$ satisfying the conditions in Theorem \ref{thm: new rigidity}.
\begin{lem}[Bounds on $G_{ij}$: small $q$ case] \label{lem: WBGij}
Let $H $ be a generalized symmetric Wigner matrix satisfying the conditions for $H$ in Theorem \ref{thm: new rigidity} (in particular $|h_{ij}|\le N^{-\phi}$ almost surely). Then, for any $0 < c < 1$, there are constants $C>0$ and $N_0$, depending only on $(C_1,\phi,c)$, such that for all $N\ge N_0$, all $z=E+i\eta$ with $|E|\le 5$ and $\eta\ge N^{-1+c}$, and all $i\ne j$, we have the following weak bound on $G_{ij}$:
\be\label{res: WBGij}
\E|G_{ij}(z)|^2 \leq \varphi^C \left(  \frac{\im m_{sc}(z)}{N \eta}  + \frac{1}{(N \eta)^2}\right).
\ee
\end{lem}


In the remainder of this section, we give the proof of Lemma \ref{lem: ex edge un}.

\begin{proof}[Proof of Lemma \ref{lem: ex edge un}]
Throughout the proof, $q:=N^{\phi}$ and $\kappa:=\phi-1/3>0$. We derive the lemma from Lemma \ref{lem: edge un str}, applied to the pairs $(H^{\V},H^{\w})$ of Lemma \ref{lem: ex edge un} with the rate $N^{-\kappa+C\e}$ in \eqref{sbc22}. Lemma \ref{lem: edge un str} is the common structure of the proofs of \cite[Theorem 2.4]{EYY} and \cite[Proposition 6.4]{EKYY2}; its proof uses neither Lemma \ref{lem: ex edge un} nor Theorem \ref{thm: new edge un}. Its hypotheses are checked as follows.

(i) The matrix $H^{\V}$ satisfies \eqref{eq:diag} with $C_0$ and, since $|h_{ij}|\le q^{-1}$ almost surely, the bounded support condition with $q$ and every $c>0$, in particular $c=1$. The GOE satisfies \eqref{eq:diag} with $C_0=1$ and, by the Gaussian tail, the bounded support condition with $N^{2/5}$ in place of $q$ and some $c>0$. For each of them, Lemma \ref{lem: old} with $\zeta=1$ gives, for every sufficiently large $C$, with high probability and with constants depending only on $(C_0,\phi)$: $\|H\|\le 2+\varphi^C(q^{-2}+N^{-2/3})\le 3$ by \eqref{sjsy}; \eqref{tby4} for all $z\in{\bf S}(C)$ with $\eta\le N^{-2/3}$, since by \eqref{scm} $|m-m_{sc}|\le\varphi^C(q^{-2}\eta^{-1/2}+(N\eta)^{-1})$ and $q^{-2}\eta^{-1/2}=(N\eta)^{-1}N^{1-2\phi}\eta^{1/2}\le(N\eta)^{-1}N^{2/3-2\phi}\le (N\eta)^{-1}$; \eqref{hh1} by the rigidity bound of Lemma \ref{lem: old} (4) at $j=N$, since $\gamma_N=2$ and $q^{-2}\le N^{-2/3}$; and \eqref{hh2} by the same rigidity bound and \eqref{def:gamma}, since $\gamma_j\ge 2-3\varphi^CN^{-2/3}$ holds for at most $C'\varphi^{3C/2}$ indices $j$. (For the GOE, $q$ is replaced by $N^{2/5}$ in these bounds.)

(ii) The comparison hypothesis \eqref{sbc22} for the pair, with test functions having four derivatives and right side $N^{-\kappa+C\e}$, is \eqref{sbc2} of the following Claim, since $N^{1/3+C\e}q^{-1}=N^{-\kappa+C\e}$.

Lemma \ref{lem: edge un str} then gives $\delta>0$ and $N_0$, depending only on $(C_0,\phi)$, such that \eqref{qset} holds for all $N\ge N_0$, all $H^{\V}$ as in Lemma \ref{lem: ex edge un} and all $s\in\mathbb R$. It remains to prove the following claim, which corresponds to Proposition 6.6 of \cite{EKYY2} with $f=0$.

\textit{Claim.} Let $C_F>0$. There exist $\widetilde{\epsilon}>0$ and $C>0$, depending only on $(C_F,C_0,\phi)$, such that for every $0<\epsilon<\widetilde\epsilon$ there is $N_0$, depending only on $(C_F,C_0,\phi,\epsilon)$, with the following property. Let $N\ge N_0$, let $H^{\V}$ be as in Lemma \ref{lem: ex edge un}, and let $F\in C^4(\mathbb R)$ satisfy
\begin{equation}\label{jhw}
\sup_x | F^{(n)} (x) | (1+ |x|)^{-C_F} \leq C_F, \quad n = 1, 2, 3, 4.
\end{equation}
Set
\begin{equation}
I_{\epsilon}:=\{ x : |x-2| \leq N^{-2/3 + \epsilon} \}, \qquad \eta := N^{-2/3 - \epsilon}.
\end{equation}
Then, for every $E\in I_\epsilon$,
\begin{equation}\label{sbc}
\Big| \E^{\V} F \left( N \eta \: \im m (E + i \eta) \right) - \E^{\w} F \left( N \eta \: \im m (E + i \eta) \right) \Big| \leq N^{1/3 + C \epsilon} q^{-1}
\end{equation}
and, for all $E_1,E_2\in I_\epsilon$ (in either order; the integral is oriented),
\begin{equation}\label{sbc2}
\left| \E^{\V} F \left( N \int_{E_1}^{E_2} dy \: \im m (y + i \eta) \right) - \E^{\w} F \left( N \int_{E_1}^{E_2} dy \: \im m (y + i \eta) \right) \right| \leq N^{1/3 + C \epsilon} q^{-1}.
\end{equation}
Only \eqref{sbc2} enters Lemma \ref{lem: edge un str}.

The differences in \eqref{sbc}--\eqref{sbc2} and below do not change if $F$ is replaced by $F-F(0)$, so we assume $F(0)=0$; then $|F(x)|\le C_F|x|(1+|x|)^{C_F}$. By \eqref{jhw}, for every fixed $a>0$ and $n\in\{1,2\}$,
$$
|F^{(n)}(x)|\le N^{(aC_F+1)\e} \quad \text{for } |x|\le 2N^{a\e} \text{ and } N\ge N_0(C_F,a,\e).
$$
For a real symmetric matrix $H$ with Green function $G$, the spectral decomposition gives the Ward identity $\sum_j|G_{ij}(z)|^2=\eta^{-1}\im G_{ii}(z)$, hence
$$
N\eta\,\im m(z)=\eta^2\sum_{i,j}|G_{ij}(z)|^2, \qquad N\int_{E_1}^{E_2}\im m(y+i\eta)\,dy=\eta\int_{E_1}^{E_2}\sum_{i,j}|G_{ij}(y+i\eta)|^2\,dy .
$$
We first prove the corresponding bounds for the off-diagonal parts:
\begin{equation} \label{eq:claim_old}
\left| \E F \left( \eta^2 \sum_{i \neq j} G^{\V}_{ij} \overline{G^{\V}_{ij}} \right) - \E F \left( \eta^2 \sum_{i \neq j} G^{\w}_{ij} \overline{G^{\w}_{ij}} \right) \right| \leq N^{1/3 + C \epsilon} q^{-1},
\end{equation}
with $G$ evaluated at $E+i\eta$, and
\begin{equation} \tag{3.20$'$}
\left| \E F \left( \eta \int_{E_1}^{E_2}\sum_{i \neq j} G^{\V}_{ij} \overline{G^{\V}_{ij}}(y+i\eta)\,dy \right) - \E F \left( \eta\int_{E_1}^{E_2} \sum_{i \neq j} G^{\w}_{ij} \overline{G^{\w}_{ij}}(y+i\eta)\,dy \right) \right| \leq N^{1/3 + C \epsilon} q^{-1}.
\end{equation}
The removed diagonal parts are estimated at the end of the proof.

\textit{Interpolation.} Fix a bijective ordering map on the index set of the independent matrix elements,
\begin{equation} \label{eq:ordering}
\Phi: \{ (k, \ell) : 1 \leq k \leq \ell \leq N \} \to \{ 1, \cdots, \gamma_{\max} = \frac{N(N+1)}{2} \},
\end{equation}
and let $H_{\gamma}$ be the Wigner matrix whose matrix elements $h_{k \ell}$ ($k\le\ell$) are independent and follow the distribution of $h_{k \ell}^{\V}$ if $\Phi(k, \ell) \leq \gamma$ and the distribution of $h_{k \ell}^{\w}$ otherwise. In particular, $H_0 = H^{\w}$ and $H_{\gamma_{\max}} = H^{\V}$. (Note that the index used here is slightly different from previous papers.) For simplicity, let
\begin{equation}
G^{\gamma}_{ij} := \left( \frac{1}{H_{\gamma} -z} \right)_{ij}.
\end{equation}
The matrices $H_{\gamma}$ and $H_{\gamma -1}$ differ only at the $(a, b)$ and $(b, a)$ elements, where $\Phi (a, b) = \gamma$. Let $v:=v_{ab} := h_{ab}^{\V}$ and $w:=w_{ab} := h_{ab}^{\w}$. We define matrices $V$ and $W$ by
\begin{equation}
V_{k \ell} := \begin{cases} (\delta_{ak} \delta_{b \ell} + \delta_{a \ell} \delta_{bk})\, v_{ab}, & a\ne b,\\ \delta_{ak}\delta_{a\ell}\, v_{aa}, & a=b,\end{cases} \qquad W_{k \ell} := \begin{cases} (\delta_{ak} \delta_{b \ell} + \delta_{a \ell} \delta_{bk})\, w_{ab}, & a\ne b,\\ \delta_{ak}\delta_{a\ell}\, w_{aa}, & a=b,\end{cases}
\end{equation}
so that we can rewrite $H_{\gamma}$ and $H_{\gamma -1}$ as
\begin{equation}
H_{\gamma} = Q + V, \quad H_{\gamma -1} = Q + W,
\end{equation}
with a matrix $Q$ satisfying $Q_{ab} = Q_{ba} = 0$. The matrix $Q$ is independent of $v$ and $w$, and its law is the same in both identities. Define Green's functions
\begin{equation} \label{eq:green_fns}
R := \frac{1}{Q-z}, \quad S := \frac{1}{H_{\gamma} - z}, \quad T := \frac{1}{H_{\gamma -1} - z}.
\end{equation}

\textit{Local law for $H_\gamma$.} Each $H_\gamma$ is a generalized Wigner matrix satisfying \eqref{eq:diag} with the constant $\max(C_0,1)$, since each of its diagonal entries comes from $H^{\V}$ or from the GOE. Let $\phi':=(1/3+\phi)/2$, so that $1/3<\phi'<\phi\le1/2$, and $q':=N^{\phi'}$. The entries of $H^{\V}$ satisfy $|h|\le N^{-\phi}\le q'^{-1}$ almost surely, and a centred Gaussian entry $h$ with $\E h^2\le 2/N$ satisfies $\p(|h|>q'^{-1})\le 2\exp(-N^{1-2\phi'}/4)$. Hence every $H_\gamma$ satisfies the bounded support condition with $q'$ and $c:=(1-2\phi')/2$ for $N\ge N_1(\phi)$, and Lemma \ref{lem: old} applies to every $H_\gamma$ with $q'$, with constants depending only on $(C_0,\phi)$ and not on $\gamma$. (In general $H_\gamma$ does not satisfy the bounded support condition with $q$ itself: for $\phi=1/2$ and a Gaussian entry $h$ of variance $1/N$, $\p(|h|>N^{-1/2})=\p(|g|>1)\approx0.32$ for a standard Gaussian $g$.)

We claim that, with high probability, uniformly in $\gamma$ and in $H^{\V}$,
\begin{equation} \label{eq:S_bound}
\max_{k, \ell} \max_{E \in I_{\epsilon}} |S_{k \ell} (E + i \eta) - \delta_{k \ell} m_{sc} (E + i \eta)| \leq N^{-1/3 + 2\epsilon}.
\end{equation}
Indeed, for $E\in I_\epsilon$ and $N$ large the point $E+i\eta$ lies in ${\bf S}(C)$, and \eqref{Gij estimate} holds simultaneously for all $z\in{\bf S}(C)$ by part (1) of Lemma \ref{lem: old}. By \eqref{esmallfake}, $\im m_{sc}(E+i\eta)\le C\sqrt{||E|-2|+\eta}\le CN^{-1/3+\e/2}$, so the right side of \eqref{Gij estimate} is at most $\varphi^C(N^{-\phi'}+CN^{-1/3+3\e/4}+N^{-1/3+\e})\le N^{-1/3+2\e}$ for $N\ge N_0(C_0,\phi,\e)$. In particular, since $|m_{sc}|\le1$ by \eqref{zmsc2}, $|S_{kk}|\le2$ and $|S_{k\ell}|\le N^{-1/3+2\e}$ for $k\ne\ell$ on this event.

Let $B:=N^{\e}q^{-1}=N^{\e-\phi}$. With high probability, uniformly in $\gamma$, all entries of $H_\gamma$ are bounded by $B$ in absolute value: almost surely for the entries from $H^{\V}$, and for the Gaussian entries because the GOE satisfies the bounded support condition with $N^{1/2-\e/2}$ in place of $q$ and $c=\e/2$ for $N\ge\lceil 8^{2/\e}\rceil$, while $N^{-1/2+\e/2}\le N^{\e-\phi}$ as $\phi\le1/2$. This bound replaces $|v_{ab}|\le q^{-1}$, which fails for Gaussian entries when $\phi=1/2$. Let $\Omega$ be the intersection of the event of \eqref{eq:S_bound} for $H_\gamma$ and the event that all entries of $H_\gamma$ are bounded by $B$; then $\p(\Omega^c)\le N^{C}e^{-c\varphi}$. Let $P:=N^{-1/3+3\e}$. We claim that, on $\Omega$,
\begin{equation} \label{eq:R_bound}
\max_{k, \ell} \max_{E \in I_{\epsilon}} |R_{k \ell} (E + i \eta) - \delta_{k \ell} m_{sc} (E + i \eta)| \leq N^{-1/3 + 3\epsilon}=P.
\end{equation}
To see \eqref{eq:R_bound}, we expand $R$ using the resolvent identity $R=S+SVR$ five times:
\begin{equation}\label{8800}
R = S + SVS + (SV)^2 S + (SV)^3 S + (SV)^4 S + (SV)^5 R.
\end{equation}
Since $V$ has at most two non-zero entries, each bounded by $B$ on $\Omega$, $|S_{k\ell}|\le2$ on $\Omega$, and $|R_{k\ell}|\le\eta^{-1}$ always, we have $|((SV)^mS)_{k\ell}|\le 2^{2m+1}B^m$ for $1\le m\le4$ and $|((SV)^5R)_{k\ell}|\le 2^{10}B^5\eta^{-1}$. Hence $|R_{k\ell}-S_{k\ell}|\le 16N^{\e-\phi}+2^{10}N^{2/3+6\e-5\phi}$, and both terms are at most $P/3$ for $N\ge N_0(\phi,\e)$, since $\phi>1/3$ and $5\phi>1+3\e$. (The five terms of \eqref{8800} suffice because $\phi>1/5+3\e/5$.) Together with \eqref{eq:S_bound} and $N^{-1/3+2\e}\le P/3$, this gives \eqref{eq:R_bound}. In particular, on $\Omega$, $|R_{kk}|\le2$ and $|R_{k\ell}|\le P$ for $k\ne\ell$. The bounds \eqref{eq:S_bound}, \eqref{eq:R_bound} and $|v|\le B$ hold on $\Omega$ simultaneously for all $E\in I_\epsilon$.

\textit{The case $a\ne b$.} Let $\Delta$ be the matrix with entries $\Delta_{k\ell}=\delta_{ak}\delta_{b\ell}+\delta_{a\ell}\delta_{bk}$, so that $V=v\Delta$. We follow the proof of Lemmas 6.7 and 6.8 of \cite{EKYY2} with $f=0$ (\cite{EKYY2}, (6.41)--(6.97)), with the following inputs and changes. The variables $v$ and $w$ are centred, $\E v^2=\E w^2=N^{-1}$, and every entry $h$ of every $H_\gamma$ satisfies
$$
\E|h|^k\le C_*N^{-1}q^{2-k}, \qquad 2\le k\le 6, \qquad C_*:=\max(1+C_0,120),
$$
which is \cite[(6.42)]{EKYY2}: for the entries from $H^{\V}$ it follows from $|h|\le q^{-1}$, and for the Gaussian entries from their moments and $q\le N^{1/2}$. The a priori bounds used in \cite{EKYY2} are \eqref{eq:S_bound}, \eqref{eq:R_bound} and $|v|\le B$ on $\Omega$. The resolvent expansion \cite[(6.58)]{EKYY2} reads $S=\sum_{m=0}^5(-v)^m(R\Delta)^mR+(-v)^6(R\Delta)^6S$, and on $\Omega$ the entries of $(R\Delta)^mR$ obey \cite[(6.59)--(6.61)]{EKYY2} with the constants $2^{2m+1}$. The functionals of \eqref{eq:claim_old} contain $S_{ij}\overline{S_{ij}}$; \cite[(6.46), (6.56), (6.66)]{EKYY2} write $S_{ij}S_{ji}$, which we read as $S_{ij}\overline{S_{ij}}$, expanding the second factor by complex conjugation. Since $v$ is real, it factors out of the coefficients
$$
K^{(s,t)}_k:=\eta^2\sum_{i\ne j,\;|\{i,j\}\cap\{a,b\}|=k}\big((R\Delta)^sR\big)_{ij}\overline{\big((R\Delta)^tR\big)_{ij}},\qquad k=0,1,2,
$$
which are functions of $Q$ alone. A Taylor expansion of $F$ to third order at $y^R:=\eta^2\sum_{i\ne j}|R_{ij}|^2$ (\cite[(6.63)--(6.76)]{EKYY2}), with the bad event $\Omega^c$ treated as in \eqref{eq:diag_bad} below, gives
$$
\Big|\E F\Big(\eta^2\sum_{i\ne j}S_{ij}\overline{S_{ij}}\Big)-\E F(y^R)-\E g(Q)\,\E v^2-\E g_3(Q)\,\E(-v)^3\Big|\le N^{-5/3+C\e}q^{-1},
$$
where
\begin{align*}
g(Q)&:=F'(y^R)\sum_{k=0,1}\big(K^{(2,0)}_k+K^{(1,1)}_k+K^{(0,2)}_k\big)+\frac12F''(y^R)\big(K^{(1,0)}_0+K^{(0,1)}_0\big)^2,\\
g_3(Q)&:=F'(y^R)\big(K^{(3,0)}_0+K^{(0,3)}_0\big).
\end{align*}
Here $g$ is the function of $Q$ in \cite[(6.75)]{EKYY2}, whose second line should read $Q^{(1)}_0$ in place of $Q^{(0)}_1$, as in its first line and in (6.74). The first-order term has expectation zero, since $Q$ is independent of $v$ and $\E v=0$. The same holds with $T$ and $w$ in place of $S$ and $v$. Since $\E v^2=\E w^2=N^{-1}$ and $Q$ has the same law in both cases, the second-order terms $\E g(Q)\,\E v^2$ and $\E g(Q)\,\E w^2$ are equal and cancel. For the third-order term, $|\E(-v)^3|\le C_*N^{-1}q^{-1}$, and $|\E g_3(Q)|\le N^{-2/3+C\e}$ by \cite[(6.77)--(6.86)]{EKYY2}: the main part of $((R\Delta)^3R)_{ij}\overline{R_{ij}}$ for $i,j\notin\{a,b\}$ is $m_{sc}^2(R_{ia}R_{bj}+R_{ib}R_{aj})\overline{R_{ij}}$, and for $i,j,a,b$ distinct, \cite[Lemma 6.8]{EKYY2} with $f=0$ gives
$$
\big|\E F'(y^R)R_{ia}R_{bj}\overline{R_{ij}}\big|\le N^{-4/3+C\e}.
$$
In \cite{EKYY2} the replaced entry is $(b,d)$, and the assumption of Lemma 6.8 there should read ``$b,d,i,j$ are all distinct''. In its proof we use \cite[Lemma 6.10]{EKYY2} for $R$ in the one-sided form $R_{ia}=-R_{aa}\sum_kQ_{ak}R^{(a)}_{ki}$ ($i\ne a$), and we replace the large deviation bound after \cite[(6.97)]{EKYY2} by the second moment bound $\E|\sum_kQ_{ak}c_k|^2\le N^{-1}\sum_k\E|c_k|^2$ for coefficients $c_k$ independent of the row $a$ of $Q$ with $c_a=0$, since that row contains both bounded and Gaussian entries. Altogether, when $a\ne b$,
\begin{equation} \label{eq:S-T 1}
\left| \E F \left( \eta^2 \sum_{i \neq j} S_{ij} \overline{S_{ij}} \right) - \E F \left( \eta^2 \sum_{i \neq j} T_{ij} \overline{T_{ij}} \right) \right| \leq N^{-5/3 + C \epsilon} q^{-1}.
\end{equation}

\textit{The case $a = b$.} Then $V=v\Delta^{(a)}$ with $v=v_{aa}$, where $\Delta^{(a)}_{k\ell}:=\delta_{ak}\delta_{a\ell}$, and the resolvent identity $S=R-RVS$, applied twice, gives
\begin{equation}
S = R - RVR + (RV)^2 S, \qquad\text{i.e.}\qquad S_{ij}=R_{ij}-vR_{ia}R_{aj}+v^2R_{ia}R_{aa}S_{aj}.
\end{equation}
Hence, for $i\ne j$,
\begin{equation} \label{eq:short_exp}
S_{ij} \overline{S_{ij}} = R_{ij} \overline{R_{ij}} - 2v\,\re\big(\overline{R_{ij}}R_{ia}R_{aj}\big) + 2v^2\,\re\big(\overline{R_{ij}}R_{ia}R_{aa}S_{aj}\big) + \big|vR_{ia}R_{aj}-v^2R_{ia}R_{aa}S_{aj}\big|^2 .
\end{equation}
On $\Omega$ we have $|v|\le B\le1/2$, $|R_{kk}|,|S_{kk}|\le2$ and $|R_{k\ell}|,|S_{k\ell}|\le P\le1$ for $k\ne\ell$. Hence, for $i\ne j$, on $\Omega$,
\begin{equation}
\begin{split}
\Big|S_{ij}\overline{S_{ij}}-R_{ij}\overline{R_{ij}}+2v\,\re\big(\overline{R_{ij}}R_{ia}R_{aj}\big)\Big|&\le\begin{cases}8v^2P^3, & a\notin\{i,j\},\\ 24v^2P^2, & a\in\{i,j\},\end{cases}\\
\big|S_{ij}\overline{S_{ij}}-R_{ij}\overline{R_{ij}}\big|&\le\begin{cases}6|v|P^3, & a\notin\{i,j\},\\ 16|v|P^2, & a\in\{i,j\}.\end{cases}
\end{split}
\end{equation}
Let
\begin{equation} \label{eq:def_y^R}
y^S := \eta^2 \sum_{i \neq j} S_{ij} \overline{S_{ij}}, \quad y^R := \eta^2 \sum_{i \neq j} R_{ij} \overline{R_{ij}}, \quad \ell(R):=-2\eta^2\sum_{i\ne j}\re\big(\overline{R_{ij}}R_{ia}R_{aj}\big).
\end{equation}
There are at most $N^2$ pairs with $a\notin\{i,j\}$ and at most $2N$ pairs with $a\in\{i,j\}$, and $NP\ge2$. Summing the previous bounds, we obtain on $\Omega$
\begin{equation}
|y^S - y^R| \leq D_1:=32B\eta^2N^2P^3, \qquad |y^S-y^R-v\,\ell(R)|\le D_2:=32B^2\eta^2N^2P^3,
\end{equation}
where $\eta^2N^2P^3=N^{-1/3+7\e}$, so that $D_1=32N^{-1/3+8\e}q^{-1}$ and $D_2=32N^{-1/3+9\e}q^{-2}$.

Consider the Taylor expansion
\begin{equation}
F(y^S) - F(y^R) = F'(y^R) (y^S - y^R) + \frac{1}{2} F'' (\zeta) (y^S - y^R)^2
\end{equation}
for some $\zeta$ between $y^S$ and $y^R$. On $\Omega$, $0\le y^S\le\eta^2N^2N^{-2/3+4\e}=N^{2\e}$ and $|y^R|\le N^{2\e}+D_1\le 2N^{2\e}$, so $|F'(y^R)|,|F''(\zeta)|\le N^{(2C_F+1)\e}$, and by the bounds on $y^S-y^R$,
$$
\big|F(y^S)-F(y^R)-F'(y^R)\,v\,\ell(R)\big|\le N^{(2C_F+1)\e}\Big(D_2+\frac{D_1^2}{2}\Big)\qquad\text{on }\Omega .
$$
The linear term $F'(y^R)\,v\,\ell(R)$ has expectation zero on the whole probability space: $y^R$ and $\ell(R)$ are bounded functions of $Q$, which is independent of $v$, and $\E v=0$. Off $\Omega$ we use the deterministic bounds $|R_{k\ell}|,|S_{k\ell}|\le\eta^{-1}\le N$, which give $|y^S|,|y^R|\le N^2$ and $|\ell(R)|\le 2N^3$, hence $|F(y^S)-F(y^R)-F'(y^R)v\ell(R)|\le\Phi_N(1+|v|)$ with $\Phi_N:=2C_FN^3(1+N^2)^{C_F}$. Since $|v|\le T+v^2/T$ for every $T>0$ and $\E v^2\le C'N^{-1}$ with $C':=1+\max(C_0,1)$, we obtain
\begin{equation} \label{eq:diag_bad}
\big|\E F(y^S)-\E F(y^R)\big|\le N^{(2C_F+1)\e}\Big(D_2+\frac{D_1^2}{2}\Big)+\Phi_N\Big((1+T)\,\p(\Omega^c)+\frac{C'}{NT}\Big).
\end{equation}
With $T:=\Phi_NN^2$ and $\p(\Omega^c)\le N^Ce^{-c\varphi}$, the last term is at most $N^{-2}\le N^{-1/3}q^{-2}$ for $N$ large. Since $D_1^2/2\le D_2$ for $N$ large, we obtain
\begin{equation}
\left| \E \big( F(y^S) - F(y^R) \big) \right| \leq N^{-1/3 + C\epsilon} q^{-2}.
\end{equation}
Only the upper bound $\E v^2\le C'N^{-1}$ was used, not the value of $\E v^2$: on the diagonal the second moments of $H^{\V}$ and of the GOE need not match (for $C_0=0$, for instance, $\E(h^{\V}_{aa})^2=N^{-1}$, while the GOE has $2N^{-1}$). The same argument applies to $H_{\gamma-1}=Q+W$, with $w\sim\mathcal N(0,2/N)$ and the event $\Omega$ for $H_{\gamma-1}$. Since $y^R$ is the same in both cases, we find that
\begin{equation} \label{eq:S-T 2}
\left| \E F \left( \eta^2 \sum_{i \neq j} S_{ij} \overline{S_{ij}} \right) - \E F \left( \eta^2 \sum_{i \neq j} T_{ij} \overline{T_{ij}} \right) \right| \leq N^{-1/3 + C \epsilon} q^{-2}.
\end{equation}

\textit{The integrated functional.} The same proof gives \eqref{eq:S-T 1} and \eqref{eq:S-T 2} for $x^S:=\eta\int_{E_1}^{E_2}\sum_{i\ne j}|S_{ij}(y+i\eta)|^2dy$ in place of $y^S$ (and likewise for $R$ and $T$), for all $E_1,E_2\in I_\epsilon$, with three changes. First, the bounds on $\Omega$ hold simultaneously for all $y\in I_\epsilon$, so the pointwise bounds above hold at every $y+i\eta$ with $y$ between $E_1$ and $E_2$; after integration in $y$ the weight $\eta^2$ is replaced by $\eta|E_2-E_1|\le 2N^{-4/3}=2N^{2\e}\eta^2$, which costs a factor $2N^{2\e}$, and on $\Omega$ the functionals are bounded by $2N^{5\e}$. Second, $v$ does not depend on $y$ and leaves the $y$-integral, so the linear term is $v$ times a function of $Q$; its expectation vanishes by Fubini's theorem, since the integrand is continuous in $(H_\gamma,y)$ and bounded by $2N^2\eta^{-3}$. Third, off $\Omega$ the deterministic bounds are multiplied by $\eta^{-1}|E_2-E_1|\le N$, which only changes $\Phi_N$. In the case $a\ne b$, the same changes apply to \cite[(6.41)--(6.97)]{EKYY2}, where the integrated functional is (6.56).

\textit{Telescoping.} We write the quantity in the left hand side of \eqref{eq:claim_old} as a telescopic sum,
\begin{equation} \label{eq:tel_1}
\E F \left( \eta^2 \sum_{i \neq j} G^{\V}_{ij} \overline{G^{\V}_{ij}} \right) - \E F \left( \eta^2 \sum_{i \neq j} G^{\w}_{ij} \overline{G^{\w}_{ij}} \right) = \sum_{\gamma=1}^{\gamma_{\max}} \left(\E F \left( \eta^2 \sum_{i \neq j} S_{ij} \overline{S_{ij}} \right) - \E F \left( \eta^2 \sum_{i \neq j} T_{ij} \overline{T_{ij}} \right)\right),
\end{equation}
where $S$ and $T$ depend on $\gamma$. Among the summands, $N$ have $a=b$ and $N(N-1)/2\le N^2$ have $a\ne b$. By \eqref{eq:S-T 1} and \eqref{eq:S-T 2}, the absolute value of \eqref{eq:tel_1} is at most
$$
N\cdot N^{-1/3+C\e}q^{-2}+N^2\cdot N^{-5/3+C\e}q^{-1}\le 2N^{1/3+C\e}q^{-1},
$$
since $N^{2/3}q^{-2}=N^{1/3}q^{-1}N^{1/3-\phi}\le N^{1/3}q^{-1}$. This gives \eqref{eq:claim_old} after replacing $C$ by $C+1$ (for $N^{\e}\ge2$), and the same telescoping sum for the integrated functional gives (3.20$'$).

\textit{The diagonal parts.} Let $H$ be $H^{\V}=H_{\gamma_{\max}}$ or $H^{\w}=H_0$, with Green function $G$, and let $\Omega_0$ be the event of \eqref{eq:S_bound} for $H$. By the Ward identity, $N\eta\,\im m(E+i\eta)-\eta^2\sum_{i\ne j}|G_{ij}|^2=\eta^2\sum_i|G_{ii}|^2$. On $\Omega_0$, $|G_{ii}|\le2$ for all $i$ and all $E\in I_\epsilon$, so this difference lies in $[0,4N\eta^2]=[0,4N^{-1/3-2\e}]$, and $0\le\eta^2\sum_{i\ne j}|G_{ij}|^2\le N^{2\e}$. Everywhere, $0\le\eta^2\sum_{i\ne j}|G_{ij}|^2\le N\eta\,\im m\le N$, so both values of $F$ are at most $C_FN(1+N)^{C_F}$ in absolute value. Hence
\begin{align*}
\Big|\E F\big(N\eta\,\im m(E+i\eta)\big)-\E F\Big(\eta^2\sum_{i\ne j}|G_{ij}|^2\Big)\Big|&\le 4N^{(2C_F+1)\e}N^{-1/3-2\e}+2C_FN(1+N)^{C_F}\,\p(\Omega_0^c)\\
&\le N^{-1/3+(2C_F+1)\e}.
\end{align*}
For the integrated functional, $N\int_{E_1}^{E_2}\im m(y+i\eta)\,dy-\eta\int_{E_1}^{E_2}\sum_{i\ne j}|G_{ij}|^2dy=\eta\int_{E_1}^{E_2}\sum_i|G_{ii}|^2dy$ is at most $4N\eta|E_2-E_1|\le 8N^{-1/3}$ in absolute value on $\Omega_0$, both functionals are bounded by $2N^{5\e}$ on $\Omega_0$ and by $2N^2$ everywhere, and the same argument gives the error $N^{-1/3+(5C_F+2)\e}$. Since $\phi\le1/2$, $N^{-1/3}\le N^{1/3}q^{-1}$. Combining these errors for $H^{\V}$ and $H^{\w}$ with \eqref{eq:claim_old} and (3.20$'$), we obtain \eqref{sbc} and \eqref{sbc2} with the bound $3N^{1/3+C\e}q^{-1}\le N^{1/3+(C+1)\e}q^{-1}$ for $N^{\e}\ge3$. This proves the claim, which implies the desired lemma.
\end{proof}

\begin{rem}[Corrections to the cited works] \label{rem: corrections}
The formalization of the results of \cite{EKYY1} and \cite{EKYY2} in the forms of Lemma \ref{lem: old} and Lemma \ref{lem: ex edge un} required the following corrections. None of them changes a statement used in this paper. The corrections to \cite{EYY} are given in the proof of Lemma \ref{lem: edge un str}.
\begin{enumerate}
\item The bound \cite[(3.22)]{EKYY1} holds as printed only for $q\le N^{1/4}$; with the Hilbert--Schmidt term $(N^{-2}\sum_{i,j}|B_{ij}|^2)^{1/2}$ added to its right side it holds in general, and this is the form used for \cite[(3.32)]{EKYY1}.
\item \cite[(3.36)]{EKYY1} and the proof of \cite[Lemma 3.15]{EKYY1} omit the term $j=i$ of $A_i$; with it, $|A_i|\le 5/(4N)+22\Lambda_o^2$.
\item In the proof of \cite[Lemma 3.16]{EKYY1}, ``$(2-1-O(\cdot))^2\ge3/2$'' should read ``$(2-1-O(\cdot))\cdot2\ge3/2$''.
\item In the proofs of \cite[Lemmas 5.7 and 5.11]{EKYY1}, the constant $1/2$ in the induction should be $1/4$.
\item In the proof of \cite[Lemma 5.13]{EKYY1}, the quantity $X$ in (5.95) needs the additional term $(N\eta)^{-1}$, and for $U=\emptyset$ the Hilbert--Schmidt term in (5.105) is controlled through $\im m^{(Si)}$.
\item \cite[(A.4)]{EKYY1} (Lemma A.1 (ii)) holds for $n\le N$ random variables and $q\ge1$, which is the case used; without these conditions it fails (for instance for $q<1$ and $n=N=2$, or for $n=N^2$ and $q=1$).
\item In \cite{EKYY2}, the display after (6.19) and the argument of $K$ in (6.17) miss a factor $\pi$: $N\int_{E_1}^{E_2}\im m(y+i\eta)\,dy$ is approximately $\pi$ times the number of eigenvalues in $[E_1,E_2]$, so $K$ is to be evaluated at this quantity divided by $\pi$.
\item In the statement of \cite[Lemma 6.8]{EKYY2}, ``$a,b,i,j$ are all distinct'' should read ``$b,d,i,j$ are all distinct''.
\item In the second line of \cite[(6.75)]{EKYY2}, $Q^{(0)}_1$ should read $Q^{(1)}_0$, as in its first line and in (6.74).
\end{enumerate}
\end{rem}



